\begin{abcliste}
\abc
Konstante Geschwindigkeit heisst keine Beschleunigung: Die resultierende Kraft ist null, der Luftwiderstand hält der Gewichtskraft das Gleichgewicht.
\begin{align}
 \ssc{F}{L} &= \FLaF \\
   &= \m\cdot\g \\
   &= \FLa \approx \result{\FLaP{0}{2}}
\end{align}

\abc
Jetzt ist der Luftwiderstand grösser als die Gewichtskraft: Die resultierende Kraft zeigt nach oben.
\begin{align}
 \ssc{F}{L} - m\,g &= m\,a \\
 a &= \aF \\
   &= \frac{\FLb-\m\cdot\g}{\m} \\
   &= \a \approx \result{\aP{0}{2}}
\end{align}
Die Beschleunigung zeigt nach oben: Sie wird mit $\agP{0}{1}\,g$ abgebremst.
\end{abcliste}
